\section{Conclusion}
\label{sec:conclusion}

We have presented FS-LoRA, a method for continual LoRA adaptation that combines
first-order adapter-subspace flatness (AS$^1$) and trace-normalized Fisher-EWC
through a decoupled gradient update.
The central finding is that these two objectives are near-orthogonal in the $(A,B)$ parameter space
($\cosff \approx 0$), a consequence of the spectral separation between the current-task Hessian
and the past-task Fisher in the effective adapter update space $\DeltaW = BA$,
mediated by the approximate isotropy of the LoRA Jacobian $JJ^\top$.

This orthogonality has a direct practical implication: the two gradient components
can be designed and tuned independently, with neither substantially interfering with the other.
FS-LoRA exploits this by keeping the GAM closure strictly on the current-task loss,
adding the Fisher gradient as an independent term, and normalizing the Fisher by its
trace to make the EWC hyperparameter dataset-agnostic.

Experiments on six fine-grained recognition datasets confirm improvements
in final average accuracy on four of six datasets relative to SeqLoRA+GAM,
with consistent reductions in forgetting.
The three-layer ablation (coupled~$\to$ decoupled~$\to$ trace-normalized)
directly validates each design decision.
On standard PECL benchmarks, LR-tuned results show FS-LoRA within $1$pp of EWC-LoRA
on ImageNet-A ($58.92$ vs.\ $59.89$), and partial runs exceed EWC-LoRA on
ImageNet-R ($+4.3$pp) and CIFAR-100 ($+0.5$pp at task~7).

\textbf{Limitations.}
FS-LoRA introduces two additional hyperparameters ($\lambdaflat$ and $\lambdaewc$) relative to SeqLoRA.
OxfordPet demonstrates that an overly large perturbation radius can degrade performance
when the task gradient scale is small.
Track~2 results (except ImageNet-A) remain partial; full comparison against EWC-LoRA
is pending run completion.
The theoretical argument for near-orthogonality rests on two regularity conditions
(spectral separation and Jacobian isotropy); the violation regime is not fully characterized,
and Jacobian isotropy is a simplifying assumption rather than a direct empirical claim.

\textbf{Future work.}
Adaptive perturbation scaling that adjusts $\rho$ or $\lambdaflat$ to the task gradient norm
would address the OxfordPet outlier and make FS-LoRA more robust across diverse task scales.
The orthogonality principle may generalize to other joint objectives in continual learning
beyond flatness and Fisher, providing a geometric criterion for principled component combination.

\begin{ack}
Use unnumbered first level headings for the acknowledgments.
Do not include this section in the anonymized submission.
\end{ack}
